\usepackage{booktabs}
%affichage correct des caractères
\usepackage[utf8]{inputenc}
\usepackage[T1]{fontenc}
\usepackage[english]{babel}%choix de la langue
\usepackage{graphicx}%images
\usepackage{subcaption}
\usepackage{float}%réglages pour les flottants
\usepackage{array}%tableaux
\usepackage{longtable}
\usepackage[colorlinks=true,linkcolor=blue, urlcolor = blue]{hyperref}
\usepackage{titlesec}
\hypersetup{breaklinks=true}
\usepackage{color}%texte en couleur
\usepackage{vmargin}%réglage de marge compatible avec fancy
\usepackage{fancyhdr}%modifie plus facilement head et foot
\usepackage[skip=5pt]{caption}%permet de personnaliser la config des légendes
\usepackage{titleref}
\usepackage{hyperref}
\usepackage{pdfpages}
\usepackage{amsmath}
\usepackage{amssymb}
\usepackage{mathrsfs}
\usepackage{tocloft}

\usepackage{booktabs}
\usepackage{multirow}
\usepackage{siunitx}
\usepackage[table]{xcolor}
\usepackage{colortbl}

% Define a custom command for code formatting
\newcommand{\code}[1]{%
  \colorbox{lightgray}{\texttt{#1}}%
}

%\usepackage{gensymb}
\usepackage{relsize,exscale}
\usepackage{xurl} % permet de faire un retour à la ligne automatique dans les urls
\usepackage{listings}
\usepackage{xcolor} % For syntax highlighting

\usepackage{algorithm}
\usepackage{algpseudocode}
\usepackage{tocloft}
\usepackage{fancyhdr}
% \usepackage[colorlinks=true, linkcolor=blue, citecolor=blue, urlcolor=blue]{hyperref}  % Hyperlink package configuration

\usepackage{listings}
\usepackage{xcolor}
\usepackage{booktabs}
\usepackage{dirtree}
\usepackage{hyperref}
\usepackage{geometry}   % optional, for margins

\definecolor{codebg}{HTML}{F5F5F5}
\definecolor{codekw}{HTML}{0000AA}
\definecolor{codecomment}{HTML}{007700}
\definecolor{codestring}{HTML}{AA0000}

\lstdefinestyle{bashstyle}{
  backgroundcolor=\color{codebg},
  basicstyle=\ttfamily\small,
  keywordstyle=\color{codekw}\bfseries,
  commentstyle=\color{codecomment}\itshape,
  stringstyle=\color{codestring},
  breaklines=true,
  frame=single,
  framesep=4pt,
  rulecolor=\color{gray!50},
  numbers=none,
  showstringspaces=false,
  tabsize=2,
}

\lstdefinestyle{pythonstyle}{
  backgroundcolor=\color{codebg},
  basicstyle=\ttfamily\small,
  language=Python,
  keywordstyle=\color{codekw}\bfseries,
  commentstyle=\color{codecomment}\itshape,
  stringstyle=\color{codestring},
  breaklines=true,
  frame=single,
  framesep=4pt,
  rulecolor=\color{gray!50},
  numbers=left,
  numberstyle=\tiny\color{gray},
  showstringspaces=false,
  tabsize=2,
}

\usepackage{nomencl}
